\par\noindent\textbf{(1)} 环上角动量算符 \(p_\varphi=-i\hbar\partial_\varphi/R\)。最小耦合 \(p\mapsto p-qA\) 得 \(H=\frac{1}{2mR^2}[-i\hbar\partial_\varphi-q\Phi/(2\pi)]^2\)。由于环上 \(\boldsymbol B=0\)，作用仍可依赖非可缩回路的磁矢势环流 \(\oint\boldsymbol A\cdot d\boldsymbol l=\Phi\)。
\par\noindent\textbf{(2)} 波函数单值要求 \(\psi_n(\varphi)=e^{in\varphi}/\sqrt{2\pi}\)、\(n\in\mathbb Z\)。对应能谱 \(E_n(\Phi)=\frac{\hbar^2}{2mR^2}(n-\Phi/\Phi_0)^2\)，其中磁通量子 \(\Phi_0=h/q\) 带电荷符号，若关心周期幅值可取 \(|h/q|\)。
\par\noindent\textbf{(3)} 单个本征态能量对磁通导数给平衡持续电流 \(I_n=-\partial E_n/\partial\Phi=\frac{\hbar^2}{mR^2\Phi_0}(n-\Phi/\Phi_0)\)。把磁通变为 \(\Phi+\Phi_0\) 时可重新标号 \(n\to n+1\)，整个能谱与平衡可观测量周期不变，而同一固定量子数的能量函数非周期。
\par\noindent\textbf{(4)} 当 \(\Phi/\Phi_0=n+1/2\)，邻近 \(n,n+1\) 两支能量相等，理想无杂质环形成简并点。规范变换可以把矢势局部消掉，却会令波函数绕环边界相位扭曲；不可同时去掉可测的AB相位 \(q\Phi/\hbar\) 和单值性要求。
